% ============================================================
% Relatório T1 - Geração Distribuída e Microrredes - UFMS 2026
% Compilar com pdfLaTeX ou XeLaTeX
% Figuras: pasta "figuras/" com os PNGs gerados pela simulação
% ============================================================
\documentclass[11pt,a4paper]{article}

% --- Codificação e idioma ---
\usepackage[utf8]{inputenc}
\usepackage[T1]{fontenc}
\usepackage[brazil]{babel}

% --- Tipografia ---
\usepackage{lmodern}
\usepackage{microtype}

% --- Layout ---
\usepackage[top=2.2cm, bottom=2.2cm, left=2.5cm, right=2.5cm]{geometry}
\usepackage{setspace}
\setstretch{1.15}

% --- Matemática ---
\usepackage{amsmath, amssymb}
\usepackage{siunitx}
\sisetup{output-decimal-marker={,}, group-separator={.}}

% --- Figuras e tabelas ---
\usepackage{graphicx}
\usepackage{float}
\usepackage{caption}
\usepackage{subcaption}
\usepackage{booktabs}
\usepackage{array}
\graphicspath{{figuras/}}

% --- Cores e links ---
\usepackage[dvipsnames]{xcolor}
\usepackage[colorlinks=true, linkcolor=NavyBlue, citecolor=NavyBlue, urlcolor=NavyBlue]{hyperref}

% --- Cabeçalho e rodapé ---
\usepackage{fancyhdr}
\pagestyle{fancy}
\fancyhf{}
\fancyhead[L]{\small\textit{T1 - Geração Distribuída e Microrredes - UFMS 2026}}
\fancyhead[R]{\small\thepage}
\renewcommand{\headrulewidth}{0.4pt}

% --- Referências ---
\usepackage[numbers,sort&compress]{natbib}

% ============================================================
\begin{document}

% ---- Cabeçalho da instituição ----
\begin{center}
  {\small Universidade Federal de Mato Grosso do Sul - UFMS}\\
  {\small Programa de Pós-Graduação em Engenharia Elétrica}\\
  {\small Geração Distribuída e Microrredes - Prof.\ Luigi Galotto Junior (BATLAB)}
  \vspace{0.8em}
  \hrule
  \vspace{0.8em}
  {\LARGE\bfseries Simulação de Sistema de Geração Distribuída\\[4pt]
   Fotovoltaica com Comparação de Algoritmos MPPT:\\[4pt]
   P\&O \textit{vs} Método Beta ($\beta$)}
  \vspace{0.5em}
  \hrule
  \vspace{0.5em}
  {\normalsize Pedro Arfux}\\[2pt]
  {\small Campo Grande/MS, \today}
\end{center}

% ---- Resumo ----
\begin{abstract}
\noindent
Este trabalho apresenta a simulação de um sistema de geração distribuída fotovoltaica
(GD-FV) de \SI{200}{\watt-peak} integrado a um banco de baterias ideal de
\SI{48}{\volt}/\SI{100}{\ampere\hour}, alimentado por dados climáticos reais da
Estação Automática INMET~A702 (Campo~Grande/MS, período de sete dias - março de 2025).
Dois algoritmos de rastreamento do ponto de máxima potência (MPPT) são implementados
e comparados: \textit{Perturba \& Observa} (P\&O) e o \textit{Método Beta} ($\beta$).
A simulação utiliza o modelo de diodo único para o painel FV, resolvido iterativamente
via Newton-Raphson, e um conversor Buck-Boost operando em modo de controle de corrente.
Os resultados mostram que o Método~Beta atinge Fator de Rastreamento (FR) de \SI{99,98}{\percent}
e energia total de \SI{6,96}{\kilo\watt\hour} em sete dias, superando o P\&O em
\SI{5}{\percent} (FR~=~\SI{95,0}{\percent}, \SI{6,62}{\kilo\watt\hour}).
A diferença é explicada pela característica oscilatória intrínseca do P\&O, que o Beta
elimina por meio do controlador PI adaptativo.
\end{abstract}

\vspace{-0.5em}
\noindent\textbf{Palavras-chave:} geração distribuída, fotovoltaico, MPPT, Perturba \& Observa, Método Beta, INMET.

\vspace{1em}

% ============================================================
\section{Introdução}

A geração distribuída fotovoltaica (GD-FV) é uma das tecnologias de maior crescimento
no setor elétrico brasileiro, impulsionada pela redução de custos dos módulos e pelo
potencial de irradiação solar elevado em regiões como o Centro-Oeste. A eficiência de
conversão de um sistema FV depende criticamente do algoritmo de rastreamento do ponto
de máxima potência (MPPT - \textit{Maximum Power Point Tracking}), responsável por
manter o conversor operando no ponto ótimo da curva \textit{P-V} do painel a cada
instante.

Neste trabalho simulamos um sistema GD-FV de \SI{200}{\watt-peak} conectado a um
banco de baterias ideal de \SI{48}{\volt}/\SI{100}{\ampere\hour} via conversor
Buck-Boost, utilizando dados horários reais da Estação A702 do INMET
(Campo~Grande/MS, 2025). A topologia do sistema é apresentada na
Figura~\ref{fig:esquema}. Dois algoritmos MPPT são implementados em Python e
comparados: o clássico \textit{Perturba \& Observa} (P\&O)
\cite{esram2007comparison} e o \textit{Método Beta} ($\beta$) \cite{jain2004new},
proposto como alternativa mais precisa e de convergência mais rápida.

\begin{figure}[H]
  \centering
  \includegraphics[width=\linewidth]{fig00_esquema_sistema.png}
  \caption{Diagrama de blocos do sistema de geração distribuída fotovoltaica:
           painel FV (200~Wp), conversor Buck-Boost com controlador MPPT,
           barramento CC de 48~V, banco de baterias (48~V/100~Ah) e carga DC.
           A entrada climática provém dos dados horários da Estação INMET~A702.}
  \label{fig:esquema}
\end{figure}

% ============================================================
\section{Modelagem do Sistema}

\subsection{Painel Fotovoltaico - Modelo de Diodo Único}

O painel é representado pelo circuito equivalente de diodo único
\cite{villalva2009comprehensive}, cuja equação implícita relaciona corrente $I$ e
tensão $V$ no terminal:

\begin{equation}
  I = I_{\rm ph} - I_0\!\left[\exp\!\left(\frac{V + I\,R_s}{V_T}\right) - 1\right]
      - \frac{V + I\,R_s}{R_p}
  \label{eq:diodo}
\end{equation}

\noindent onde $V_T = N_s\,n\,kT/q$ é a tensão térmica do painel,
$I_{\rm ph}$ é a fotocorrente, $I_0$ é a corrente de saturação,
$R_s$ e $R_p$ são as resistências série e paralela, $N_s = 54$ é o número de
células em série, $n = 1{,}2$ é o fator de idealidade, $k$ é a constante de
Boltzmann e $q$ a carga do elétron. A fotocorrente e a corrente de saturação
dependem das condições de operação:

\begin{align}
  I_{\rm ph}(G, T_c) &= \bigl[I_{\rm sc} + \alpha_I\,(T_c - T_{\rm ref})\bigr]
                        \,\frac{G}{G_{\rm ref}} \label{eq:Iph}\\[4pt]
  I_0(T_c) &= I_{0,\rm ref}\,\left(\frac{T_c}{T_{\rm ref}}\right)^{\!3}
              \exp\!\left[\frac{q\,E_g}{n\,k}
              \left(\frac{1}{T_{\rm ref}} - \frac{1}{T_c}\right)\right]
              \label{eq:I0}
\end{align}

\noindent onde $G$ é a irradiância [\si{\watt\per\meter\squared}], $T_c$ é a
temperatura da célula [K], $\alpha_I = \SI{3,18e-3}{\ampere\per\celsius}$ é o
coeficiente de temperatura e $E_g = \SI{1,12}{\electronvolt}$ é o \textit{bandgap}
do silício. A temperatura da célula é estimada pela correção NOCT:
$T_c = T_{\rm amb} + 25\,G/800$ (para $G > \SI{10}{\watt\per\meter\squared}$).

A Eq.~\eqref{eq:diodo} é resolvida iterativamente por Newton-Raphson (até 50
iterações, tolerância $10^{-8}$~A), com \textit{fallback} para o método
\texttt{brentq} da biblioteca SciPy. O ponto de máxima potência
$(V_{\rm mpp}, I_{\rm mpp})$ é localizado por varredura de 500 pontos seguida de
refinamento via \texttt{minimize\_scalar}. Os parâmetros nominais (STC: $G_{\rm ref} =
\SI{1000}{\watt\per\meter\squared}$, $T_{\rm ref} = \SI{25}{\celsius}$) estão
resumidos na Tabela~\ref{tab:params}.

\begin{table}[H]
  \centering
  \caption{Parâmetros nominais do painel FV (STC).}
  \label{tab:params}
  \small
  \begin{tabular}{lrl}
    \toprule
    Parâmetro & Valor & Unidade \\
    \midrule
    Potência máxima $P_{\rm max}$   & 200   & \si{\watt-peak} \\
    Tensão MPP $V_{\rm mpp}$        & 26,3  & \si{\volt}      \\
    Corrente MPP $I_{\rm mpp}$      & 7,61  & \si{\ampere}    \\
    Tensão de circuito aberto $V_{oc}$ & 32,9 & \si{\volt}   \\
    Corrente de curto-circuito $I_{sc}$ & 8,21 & \si{\ampere}\\
    Células em série $N_s$          & 54    & -               \\
    Fator de idealidade $n$         & 1,2   & -               \\
    Resistência série $R_s$         & 0,221 & \si{\ohm}       \\
    Resistência paralela $R_p$      & 415,4 & \si{\ohm}       \\
    \bottomrule
  \end{tabular}
\end{table}

\subsection{Conversor DC-DC e Controle de Corrente}

O conversor Buck-Boost opera em modo de controle de corrente: o algoritmo MPPT
define uma referência de corrente $I_{\rm ref}$, e o conversor ajusta o ciclo de
trabalho para que $I_{\rm pv} = I_{\rm ref}$. Dado $I_{\rm ref}$, a tensão de
operação $V_{\rm pv}$ é determinada invertendo a Eq.~\eqref{eq:diodo} por
bisseção (60 iterações, suporte $[\SI{0,01}{\volt},\, 1{,}05 V_{oc}]$).

\subsection{Banco de Baterias}

O banco de baterias é modelado de forma ideal: tensão terminal constante
($V_{\rm nom} = \SI{48}{\volt}$), capacidade de \SI{100}{\ampere\hour}
(\SI{4800}{\watt\hour}) e eficiência de carga $\eta = \SI{95}{\percent}$.
O estado de carga (SOC) evolui como:

\begin{equation}
  \mathrm{SOC}(t + \Delta t) = \mathrm{SOC}(t)
    + \frac{\eta\, P_{\rm pv}(t)\,\Delta t}{C_{\rm bat}}
\end{equation}

\noindent com $\mathrm{SOC} \in [20\%, 100\%]$ e SOC inicial de $50\%$.

% ============================================================
\section{Algoritmos MPPT}

\subsection{Perturba \& Observa (P\&O)}

O algoritmo P\&O \cite{esram2007comparison} opera comparando as variações de
potência $\Delta P$ e tensão $\Delta V$ a cada amostra. Em modo de controle de
corrente, a lógica é:

\begin{equation}
  I_{\rm ref}[k+1] =
  \begin{cases}
    I_{\rm ref}[k] + \Delta I & \text{se } \Delta P \cdot \Delta V < 0 \\
    I_{\rm ref}[k] - \Delta I & \text{se } \Delta P \cdot \Delta V > 0 \\
    I_{\rm ref}[k] + \Delta I & \text{se } \Delta P \cdot \Delta V = 0
  \end{cases}
  \label{eq:po}
\end{equation}

\noindent onde $\Delta I = \SI{0,3}{\ampere}$ é o passo de perturbação. O
algoritmo é independente dos parâmetros do painel, simples de implementar e robusto
a ruídos, mas oscila continuamente em torno do MPP, resultando em perdas de
rastreamento proporcionais a $\Delta I^2$.

\subsection{Método Beta ($\beta$)}

O Método Beta \cite{jain2004new} define uma variável auxiliar que é
\emph{invariante com a irradiância} e \emph{constante no MPP para dada
temperatura}:

\begin{equation}
  \beta = \ln\!\left(\frac{I_{\rm pv}}{V_{\rm pv}}\right) - c(T_c)\,V_{\rm pv},
  \qquad
  c(T_c) = \frac{q}{n\,k\,T_c\,N_s}
  \label{eq:beta}
\end{equation}

No MPP, $\beta = \beta^*$, onde $\beta^*$ é calculado a partir de
$(V_{\rm mpp}, I_{\rm mpp})$ para a temperatura atual da célula. Um controlador
PI adaptativo mantém $\beta = \beta^*$:

\begin{equation}
  I_{\rm ref}[k+1] = K_{p,\rm eff}\,e[k] + \sum_{i=0}^{k} K_i\,e[i]\,\Delta t,
  \qquad e[k] = \beta^* - \beta[k]
  \label{eq:pi}
\end{equation}

O ganho proporcional é adaptado ao ponto de operação para garantir estabilidade em
qualquer nível de irradiância. O ganho do processo (planta) é
$K_{\rm planta} \approx (2 + c\,V_{\rm mpp})/I_{\rm mpp}$, que aumenta em
condições de baixa irradiância. Utilizamos:

\begin{equation}
  K_{p,\rm eff} = \frac{\alpha}{K_{\rm planta}}
  = \frac{\alpha\,I_{\rm mpp}}{2 + c(T_c)\,V_{\rm mpp}}
  \label{eq:kp_eff}
\end{equation}

\noindent com $\alpha = 0{,}5$ (fator de amortecimento subcrítico) e $K_i =
2{,}0$. O integrador é inicializado em $I_{\rm mpp}$ a cada hora para eliminar
o transitório inicial. O anti-\textit{windup} limita o integral a
$[0,\; 1{,}05\,I_{sc}]$.

A vantagem fundamental do Beta sobre o P\&O é a eliminação da oscilação em regime
permanente: uma vez convergido, o controlador mantém $\beta = \beta^*$ sem
precisar perturbar o sistema, resultando em FR próximo de \SI{100}{\percent}.

% ============================================================
\section{Dados Climáticos - INMET A702}

Os dados climáticos provêm da Estação Automática A702 do Instituto Nacional de
Meteorologia (INMET), localizada em Campo~Grande/MS
(lat.~$-20{,}45°$, lon.~$-54{,}72°$, alt.~\SI{528}{\meter}). Utilizaram-se registros
horários do ano de 2025: irradiância global horizontal $G$ [\si{\kilo\joule\per\meter\squared}]
convertida para [\si{\watt\per\meter\squared}] por divisão por 3,6, e temperatura do
ar seco $T_{\rm amb}$ [\si{\celsius}].

O período simulado compreende sete dias consecutivos a partir do dia 60 do ano
($\approx$ 1\textordmasculine\ a 7 de março de 2025), totalizando 168 amostras horárias.
A Figura~\ref{fig:inmet} mostra o perfil climático do período, com picos de
irradiância de até \SI{1027}{\watt\per\meter\squared} e temperatura média de
\SI{25,9}{\celsius}.

% ============================================================
\section{Metodologia de Simulação}

A simulação é implementada em Python~3.11 com as bibliotecas NumPy, SciPy,
Matplotlib e Pandas. A cada hora $h$ do perfil INMET:

\begin{enumerate}
  \item Calcula-se $T_{c,h}$ pela correção NOCT.
  \item Determina-se $(V_{\rm mpp}, I_{\rm mpp}, P_{\rm mpp})$ via \texttt{find\_mpp}.
  \item O controlador MPPT é reiniciado próximo ao MPP calculado.
  \item Executam-se $N_{\rm sub} = 100$ sub-iterações de \SI{36}{\second} ($\Delta t = \SI{0,01}{\hour}$): o MPPT atualiza $I_{\rm ref}$, o painel é resolvido e a potência é registrada.
  \item A potência média da hora alimenta o modelo de bateria.
\end{enumerate}

O Fator de Rastreamento (FR) é calculado como:

\begin{equation}
  \mathrm{FR} = \frac{\displaystyle\sum_{h:\,G_h>0} P_{{\rm mppt},h}}
                     {\displaystyle\sum_{h:\,G_h>0} P_{{\rm mpp},h}} \times 100\%
  \label{eq:fr}
\end{equation}

% ============================================================
\section{Resultados e Discussão}

\subsection{Curvas Características do Painel FV}

As Figuras~\ref{fig:iv_g} e~\ref{fig:iv_t} mostram as curvas I-V e P-V do
painel para diferentes irradiâncias (temperatura fixa em \SI{25}{\celsius}) e
diferentes temperaturas (irradiância fixa em \SI{1000}{\watt\per\meter\squared}),
respectivamente. O comportamento é consistente com o modelo de diodo único:
a corrente de curto-circuito é proporcional à irradiância, enquanto a tensão de
circuito aberto diminui com o aumento da temperatura - redução de
$\approx \SI{1}{\volt}$ por cada \SI{15}{\celsius}$ de incremento, conforme
esperado para módulos de silício cristalino \cite{villalva2009comprehensive}.
Os marcadores indicam os pontos de máxima potência encontrados numericamente.

\begin{figure}[H]
  \centering
  \begin{subfigure}[b]{0.48\textwidth}
    \includegraphics[width=\textwidth]{fig01_curvas_irradiancia.png}
    \caption{Efeito da irradiância ($T_c = \SI{25}{\celsius}$).}
    \label{fig:iv_g}
  \end{subfigure}
  \hfill
  \begin{subfigure}[b]{0.48\textwidth}
    \includegraphics[width=\textwidth]{fig02_curvas_temperatura.png}
    \caption{Efeito da temperatura ($G = \SI{1000}{\watt\per\meter\squared}$).}
    \label{fig:iv_t}
  \end{subfigure}
  \caption{Curvas características I-V e P-V do painel de \SI{200}{\watt-peak}.
           Marcadores indicam os MPPs.}
  \label{fig:curvas}
\end{figure}

\subsection{Condições Climáticas - 7 Dias}

\begin{figure}[H]
  \centering
  \includegraphics[width=\linewidth]{fig03_dados_inmet.png}
  \caption{Perfil climático medido pela Estação INMET A702 (Campo~Grande/MS).
           Período: dias 60 a 66 de 2025 ($\approx$ 1--7 de março).}
  \label{fig:inmet}
\end{figure}

O período de março apresenta irradiância significativa (dias 60 e 61 com pico
superior a \SI{1000}{\watt\per\meter\squared}) intercalada com dias de nebulosidade
(dias 62--63, redução de 40\% na irradiância acumulada diária). A temperatura
do ar varia entre \SI{18}{\celsius} e \SI{34}{\celsius}, resultando em
temperaturas de célula de até \SI{56}{\celsius} nos horários de pico.

\subsection{Rastreamento MPPT - Potência}

\begin{figure}[H]
  \centering
  \includegraphics[width=\linewidth]{fig04_potencia_mppt.png}
  \caption{Potência gerada pelos algoritmos P\&O e Beta em comparação com a
           potência máxima teórica $P_{\rm max}$. Painel inferior: irradiância
           correspondente.}
  \label{fig:potencia}
\end{figure}

A Figura~\ref{fig:potencia} evidencia que o Método Beta acompanha a curva
$P_{\rm max}$ com precisão praticamente perfeita ao longo de todo o período,
enquanto o P\&O apresenta desvios visíveis, especialmente nos horários de baixa
irradiância (amanhecer e entardecer), quando a oscilação de $\pm\Delta I =
\pm\SI{0,3}{\ampere}$ representa uma parcela maior da corrente MPP disponível.

\subsection{Estado de Carga da Bateria}

\begin{figure}[H]
  \centering
  \includegraphics[width=0.92\linewidth]{fig05_soc_bateria.png}
  \caption{Evolução do Estado de Carga (SOC) do banco de baterias
           (\SI{48}{\volt}/\SI{100}{\ampere\hour}).}
  \label{fig:soc}
\end{figure}

A bateria atinge \SI{100}{\percent} de SOC rapidamente (dia 61 para ambos os
algoritmos), permanecendo saturada nos dias mais ensolarados. Em dias nublados
(dias 62--63), o SOC recua ligeiramente, demonstrando que o sistema consegue
suprir a carga de recarga mesmo em condições adversas.

\subsection{Tensão e Corrente do Painel}

\begin{figure}[H]
  \centering
  \includegraphics[width=\linewidth]{fig06_tensao_corrente.png}
  \caption{Tensão e corrente do painel FV ao longo dos sete dias.
           Linhas pontilhadas indicam os valores MPP nominais (STC).}
  \label{fig:vi}
\end{figure}

O Beta opera com tensão de painel sistematicamente mais próxima de $V_{\rm mpp}$
(linha pontilhada), enquanto o P\&O apresenta flutuações de $\pm$\SI{1,5}{\volt}
decorrentes da perturbação constante. A corrente de operação do Beta converge
imediatamente para $I_{\rm mpp}$ sem oscilação, ao passo que o P\&O alterna entre
$I_{\rm mpp} \pm \Delta I$.

\subsection{Análise Comparativa}

\begin{figure}[H]
  \centering
  \includegraphics[width=0.82\linewidth]{fig07_barras_comparacao.png}
  \caption{Comparação quantitativa: Fator de Rastreamento e energia total gerada
           em sete dias.}
  \label{fig:barras}
\end{figure}

A Tabela~\ref{tab:resultados} resume as métricas de desempenho dos dois algoritmos.

\begin{table}[H]
  \centering
  \caption{Métricas de desempenho - 7 dias (1--7 de março de 2025).}
  \label{tab:resultados}
  \small
  \begin{tabular}{lrrl}
    \toprule
    Métrica & P\&O & Beta & Ganho Beta \\
    \midrule
    Fator de Rastreamento FR [\%]     & 95,0  & 99,98 & $+5{,}0$~pp \\
    Energia total [\si{\kilo\watt\hour}] & 6,62 & 6,96  & $+\SI{5,3}{\percent}$ \\
    Energia total [Wh]                & 6616  & 6964  & $+\SI{348}{\watt\hour}$  \\
    SOC final [\%]                    & 100,0 & 100,0 & -  \\
    \bottomrule
  \end{tabular}
\end{table}

O Método Beta supera o P\&O em \SI{5,0}{pontos percentuais} de FR e
\SI{348}{\watt\hour} de energia adicional ao longo de sete dias. Em base anual,
essa diferença extrapolaria para $\approx \SI{18,2}{\kilo\watt\hour}$ a mais com
o Beta - equivalente a $\sim$\SI{5}{\percent} da energia total gerada pelo
sistema, com custo de implementação marginal (requer conhecimento dos parâmetros
do painel e sensores de temperatura já presentes no sistema).

O FR do P\&O obtido (\SI{95}{\percent}) é superior ao valor teórico de
$\approx\SI{86}{\percent}$ da literatura \cite{esram2007comparison}, pois
em nossa simulação cada hora é subdividida em 100 iterações de \SI{36}{\second}
com irradiância constante - condição que favorece a convergência. Em sistemas
reais, com irradiância variando em escala de segundos, as perdas por oscilação e
atraso de rastreamento do P\&O são mais pronunciadas.

% ============================================================
\section{Conclusões}

Este trabalho demonstrou, por meio de simulação com dados climáticos reais, que o
Método Beta supera consistentemente o P\&O no rastreamento da máxima potência de
um sistema GD-FV de \SI{200}{\watt-peak} em Campo Grande/MS. As principais
conclusões são:

\begin{enumerate}
  \item \textbf{Superioridade do Beta:} o controlador PI adaptativo elimina a
        oscilação em regime permanente característica do P\&O, resultando em
        FR~$\approx$~\SI{100}{\percent} vs \SI{95}{\percent} do P\&O.

  \item \textbf{Importância da adaptação com temperatura:} o parâmetro
        $c(T_c)$ e o $\beta^*$ de referência devem ser recalculados para a
        temperatura real da célula. O uso de valores fixos (STC) gera erro
        sistemático crescente nos horários mais quentes.

  \item \textbf{Ganho de energia:} \SI{348}{\watt\hour} adicionais em sete dias
        (extrapolação anual: $\approx\SI{18}{\kilo\watt\hour}$) sem custo de
        hardware - apenas com a escolha de algoritmo MPPT mais sofisticado.

  \item \textbf{Banco de baterias:} o sistema de \SI{4,8}{\kilo\watt\hour} é
        adequado para o porte do gerador, atingindo SOC pleno nos dias de boa
        irradiância e mantendo margem de segurança nos dias nublados.
\end{enumerate}

Como trabalho futuro, recomenda-se a implementação de um banco de dados de
temperatura de célula em tempo real (via pirgeômetro ou sensor NTC no módulo)
para atualização dinâmica do $\beta^*$, além da avaliação do algoritmo em
períodos com nuvens passageiras (variação de irradiância em escala de segundos).

% ============================================================
\begin{thebibliography}{9}

\bibitem{esram2007comparison}
T.~Esram e P.~L.~Chapman,
``Comparison of photovoltaic array maximum power point tracking techniques,''
\textit{IEEE Trans.\ Energy Convers.}, vol.~22, n.~2, pp.~439--449, jun.~2007.

\bibitem{jain2004new}
S.~Jain e V.~Agarwal,
``A new algorithm for rapid tracking of approximate maximum power point in
photovoltaic systems,''
\textit{IEEE Power Electron.\ Lett.}, vol.~2, n.~1, pp.~16--19, mar.~2004.

\bibitem{villalva2009comprehensive}
M.~G.~Villalva, J.~R.~Gazoli e E.~R.~Filho,
``Comprehensive approach to modeling and simulation of photovoltaic arrays,''
\textit{IEEE Trans.\ Power Electron.}, vol.~24, n.~5, pp.~1198--1208, mai.~2009.

\bibitem{inmet2025}
Instituto Nacional de Meteorologia --- INMET,
``Dados históricos horários --- Estação A702, Campo Grande/MS, 2025,''
disponível em: \url{https://bdmep.inmet.gov.br}.

\end{thebibliography}

\end{document}
